\section{Ownership sheets, births, deaths and exchanges}
\label{sec:party_sheets}

An ownership sheet assigns an owner to every site of a finite set $\Lambda$, and the register
$\C^q$ of a site is held by its owner. This section implements the homogeneous births and
deaths of \cite[Lemma~6.5]{OpenAI2026PolynomialPEPS} and the two-sheet exchange of
\cite[Lemma~6.6]{OpenAI2026PolynomialPEPS} by words and monomials of the party calculus, for
given splitting data, and assembles the memory of several sheets used in
\cite[Proposition~7.1]{OpenAI2026PolynomialPEPS}: lifts of changes of one or two sheets to the
whole stack, bounded changes, and the private product initialization. Sheets carry no encoded
holes, so the hole parts of the sets of the two lemmas are empty.

\subsection{Framing, relabelling and sequencing of original circuits}

\begin{definition}[Framing, relabelling and sequencing of original circuits]
    \label{def:party_original_frame_sequence}
    \lean{QuantumCircuit.Party.prodVec, QuantumCircuit.Party.prodVec_apply,
        QuantumCircuit.Party.norm_prodVec, QuantumCircuit.Party.prodVec_sub_left,
        QuantumCircuit.Party.prodVec_sub_right, QuantumCircuit.Party.OriginalGate.frame,
        QuantumCircuit.Party.OriginalGate.toWeightedGate_frame,
        QuantumCircuit.Party.OriginalCircuit.frame,
        QuantumCircuit.Party.OriginalCircuit.toEffectCircuit_frame,
        QuantumCircuit.Party.OriginalCircuit.nonprivateCount_frame,
        QuantumCircuit.Party.OriginalCircuit.participation_frame,
        QuantumCircuit.Party.OriginalCircuit.frameLocationsEquiv,
        QuantumCircuit.Party.OriginalCircuit.participants_frameLocationsEquiv,
        QuantumCircuit.Party.OriginalCircuit.forall_participants_frame,
        QuantumCircuit.Party.OriginalCircuit.isExpansionBounded_frame,
        QuantumCircuit.Party.OriginalCircuit.isMonomialBounded_frame,
        QuantumCircuit.Party.OriginalCircuit.isSourceBounded_frame,
        QuantumCircuit.Party.EffectCircuit.isExpansionBounded_frame,
        QuantumCircuit.Party.EffectCircuit.isMonomialBounded_frame,
        QuantumCircuit.Party.EffectCircuit.isSourceBounded_frame,
        QuantumCircuit.Party.OriginalCircuit.relabel,
        QuantumCircuit.Party.OriginalCircuit.toPartyCircuit_relabel,
        QuantumCircuit.Party.OriginalCircuit.nonprivateCount_relabel,
        QuantumCircuit.Party.OriginalCircuit.participation_relabel,
        QuantumCircuit.Party.OriginalCircuit.sequence,
        QuantumCircuit.Party.OriginalCircuit.sequence_zero,
        QuantumCircuit.Party.OriginalCircuit.sequence_succ,
        QuantumCircuit.Party.OriginalCircuit.op_sequence_zero,
        QuantumCircuit.Party.OriginalCircuit.op_sequence_succ,
        QuantumCircuit.Party.OriginalCircuit.isExpansionBounded_sequence,
        QuantumCircuit.Party.OriginalCircuit.isMonomialBounded_sequence,
        QuantumCircuit.Party.OriginalCircuit.isSourceBounded_sequence,
        QuantumCircuit.Party.OriginalCircuit.nonprivateCount_sequence,
        QuantumCircuit.Party.OriginalCircuit.forall_participants_sequence,
        QuantumCircuit.Party.OriginalCircuit.participation_sequence}
    \leanok
    \uses{def:party_original_circuit, def:party_gate_frame, def:party_relabel,
        def:party_gate_counts, def:party_weighted_gate, lem:kronecker_vectors}
    Let $X$, $Y$ and $W$ be systems on the parties. In this section $\operatorname{op}(w)$
    denotes the operator of an original circuit, word or monomial $w$, the letter $T$ being kept
    for a region of sites. For an original circuit, its number of
    nonprivate gates is its number of gates, every gate costing $1$ and local contractions
    nothing, and the participation of a party is the number of gates whose cell contains it.
    It satisfies the expansion bound $(r,S)$ if every monomial of every gate expansion has at
    most $r$ pair effects and the absolute values of the coefficients of every gate expansion
    sum to at most $S$, the monomial bound $K$ if every gate expansion has at most $K$
    monomials, and the pair-source bound $r$ if every monomial of every gate expansion has at
    most $r$ pair sources \cite[Theorem~5.2]{OpenAI2026PolynomialPEPS}.

    \emph{Product vectors.} The product system $W\times X$ has the basis $W_p\times X_p$ at the
    party $p$, and its configurations are the pairs of configurations of $W$ and $X$. For
    $x\in\mathcal H_W$ and $y\in\mathcal H_X$ the product vector
    $x\otimes y\in\mathcal H_{W\times X}$ has the coefficient $x(\sigma)y(\tau)$ at the
    configuration paired with $(\sigma,\tau)$. Thus
    $\lVert x\otimes y\rVert=\lVert x\rVert\,\lVert y\rVert$,
    $x\otimes y-x'\otimes y=(x-x')\otimes y$ and $x\otimes y-x\otimes y'=x\otimes(y-y')$.

    \emph{Framing.} For an original circuit $w$ from $X$ to $Y$, the framed circuit $w^W$ from
    $W\times X$ to $W\times Y$ runs $w$ beside the spectator factor $W_p$ placed in front of the
    basis of every party $p$: every original gate is replaced by its framed gate expansion, and
    every local contraction $A$ at a party $p$ by $1_{W_p}\otimes A$. The effect circuit of
    $w^W$ is the framed effect circuit of $w$. Framing keeps the number of nonprivate gates,
    the participation of every party, the gate occurrences together with their cells, hence
    every condition on the participants of the gates, and the expansion, monomial and
    pair-source bounds, of original circuits and of effect circuits alike.

    \emph{Relabelling.} For bijections $e_p:Y_p\to X_p$, the relabelling along $e$ is the
    original circuit from $X$ to $Y$ made of one local operation at every party, the matrix of
    $e_p$; it is the relabelling of party circuits, has no nonprivate gate, and no party
    participates in it.

    \emph{Sequencing.} For systems $a_0,a_1,\dots$ and original circuits $c_t$ from $a_t$ to
    $a_{t+1}$, the sequence of length $t$ is the empty circuit for $t=0$ and $c_0$ followed by
    $c_1,\dots,c_{t-1}$ in general, so the sequence of length $t+1$ is that of length $t$
    followed by $c_t$. Its operator is $\operatorname{op}(c_{t-1})\cdots\operatorname{op}(c_0)$,
    the identity for $t=0$. Its number of nonprivate gates and the participation of every party
    are the sums of those of the stages $c_s$, $s<t$; an expansion, monomial or pair-source
    bound satisfied by every stage $c_s$ with $s<t$ is satisfied by the sequence, and so is
    every condition on the participants of the gates.
\end{definition}

\begin{theorem}[Framed circuits and relabellings on product vectors]
    \label{thm:party_frame_prod_vec}
    \lean{QuantumCircuit.Party.OriginalCircuit.op_frame,
        QuantumCircuit.Party.OriginalGate.isCompatible_frame,
        QuantumCircuit.Party.OriginalCircuit.toEuclideanLin_op_frame_prodVec,
        QuantumCircuit.Party.OriginalCircuit.norm_op_frame_prodVec_sub_le,
        QuantumCircuit.Party.OriginalCircuit.op_relabel,
        QuantumCircuit.Party.OriginalCircuit.toEuclideanLin_op_relabel,
        QuantumCircuit.Party.OriginalCircuit.toEuclideanLin_op_relabel_prodComm_prodVec}
    \leanok
    \uses{def:party_original_frame_sequence}
    The operator of a framed original gate is the identity on the spectator factors tensored
    with the operator of the gate. For an original circuit $w$ from $X$ to $Y$ and a system $W$,
    read along the pairing of configurations,
    $\operatorname{op}(w^W)=1_{\mathcal H_W}\otimes\operatorname{op}(w)$. Hence
    $\operatorname{op}(w^W)(x\otimes y)=x\otimes\operatorname{op}(w)y$ for
    $x\in\mathcal H_W$ and $y\in\mathcal H_X$, and if $\lVert x\rVert\leq1$ then
    \begin{align}
        \lVert\operatorname{op}(w^W)(x\otimes v)-x\otimes v'\rVert
        \leq\lVert\operatorname{op}(w)v-v'\rVert
    \end{align}
    for all $v\in\mathcal H_X$ and $v'\in\mathcal H_Y$. The operator of the relabelling along
    bijections $e_p:Y_p\to X_p$ is the matrix of the bijection $\tau\mapsto(e_p(\tau_p))_p$ of
    configurations; it maps a vector $v\in\mathcal H_X$ to the vector
    $\tau\mapsto v((e_p(\tau_p))_p)$ of $\mathcal H_Y$. In particular the relabelling exchanging
    the two factors of $X\times Y$ maps $x\otimes y$ to $y\otimes x$.
\end{theorem}
\begin{proof}\leanok
    \uses{thm:party_op_frame, def:party_relabel, def:party_weighted_gate}
    The compatibility of framed original gates is that of their framed gate expansions, and
    Theorem~\ref{thm:party_op_frame} for party circuits then gives
    $\operatorname{op}(w^W)=1\otimes\operatorname{op}(w)$. Applied to $x\otimes y$, the sum over
    the first component reduces to its diagonal term, which gives
    $x\otimes\operatorname{op}(w)y$. The bound follows from
    $\operatorname{op}(w^W)(x\otimes v)-x\otimes v'=x\otimes(\operatorname{op}(w)v-v')$ and
    $\lVert x\otimes u\rVert=\lVert x\rVert\,\lVert u\rVert$. The operator of a relabelling is that
    of the relabelling of party circuits, a permutation matrix, and its action on vectors and on
    product vectors is a direct computation.
\end{proof}

\subsection{Ownership sheets}

Throughout, $\Lambda$ is a finite set of sites, $q\in\N$, $[q]=\{0,\dots,q-1\}$, and
$[q]^R$ is the set of configurations of a region $R\subseteq\Lambda$. The parties form a set
$\mathcal P$.

\begin{definition}[Ownership sheets and their reference vectors]
    \label{def:party_sheet}
    \lean{QuantumCircuit.Party.Sheet, QuantumCircuit.Party.Sheet.system,
        QuantumCircuit.Party.Sheet.refVec, QuantumCircuit.Party.Sheet.refVec_apply,
        QuantumCircuit.Party.Sheet.norm_refVec, QuantumCircuit.Party.Sheet.norm_refVec_le,
        QuantumCircuit.Party.Sheet.twoConfigEquiv,
        QuantumCircuit.Party.Sheet.twoConfigEquiv_apply,
        QuantumCircuit.Party.Sheet.twoConfigEquiv_fst_eq,
        QuantumCircuit.Party.Sheet.twoConfigEquiv_snd_eq,
        QuantumCircuit.Party.Sheet.prodVec_refVec}
    \leanok
    \uses{def:sites_block, def:party_original_frame_sequence, lem:kronecker_vectors}
    A sheet $S$ on $\Lambda$ is an owner map $o_S:\Lambda\to\mathcal P$; the register $\C^q$ of
    a site $x$ is held by $o_S(x)$ \cite[Section~6]{OpenAI2026PolynomialPEPS}. Its system $X^S$
    has at the party $p$ the basis $X^S_p=[q]^{o_S^{-1}(p)}$ of configurations of the sites $p$
    owns, the blocked system of $[q]^\Lambda$ along $o_S$. Its canonical coordinates are the
    bijection $\gamma_S$ from the configurations of $X^S$ to $[q]^\Lambda$ with
    $\gamma_S(z)(x)=z_{o_S(x)}(x)$. A vector $\Omega\in\C^{[q]^\Lambda}$ read on the sheet is
    $\Omega_S=\Omega\circ\gamma_S$, and $\lVert\Omega_S\rVert=\lVert\Omega\rVert$, so
    $\lVert\Omega_S\rVert\leq1$ when $\lVert\Omega\rVert\leq1$.

    Two sheets $S_1,S_2$ held at once form the product system $X^{S_1}\times X^{S_2}$. Its
    canonical coordinates are the bijection $\gamma_{S_1,S_2}$ onto $[q]^\Lambda\times[q]^\Lambda$,
    $\gamma_{S_1,S_2}(z)=(\gamma_{S_1}(z^1),\gamma_{S_2}(z^2))$ for the pair $(z^1,z^2)$ of
    configurations paired with $z$: the coordinate of a site $x$ on the first sheet is read at
    $o_{S_1}(x)$, on the second at $o_{S_2}(x)$. The product of the reference vectors is
    $\Omega_{S_1}\otimes\Omega_{S_2}=(\Omega\otimes\Omega)\circ\gamma_{S_1,S_2}$.
\end{definition}

\begin{definition}[Change of owner and the sets of a birth]
    \label{def:party_change_owner}
    \lean{QuantumCircuit.Party.Sheet.changeOwner, QuantumCircuit.Party.Sheet.birthEnv,
        QuantumCircuit.Party.Sheet.mem_birthEnv, QuantumCircuit.Party.Sheet.disjoint_birthEnv,
        QuantumCircuit.Party.Sheet.changeOwner_eq_self,
        QuantumCircuit.Party.Sheet.changeOwner_birth_owners,
        QuantumCircuit.Party.Sheet.owner_eq_iff_of_birth,
        QuantumCircuit.Party.Sheet.changeOwner_owner_eq_self_iff,
        QuantumCircuit.Party.Sheet.changeOwner_owner_eq_iff,
        QuantumCircuit.Party.Sheet.changeOwner_owner_eq_iff_of_ne,
        QuantumCircuit.Party.Sheet.changeOwner_self_owner_eq_iff,
        QuantumCircuit.Party.Sheet.mem_sdiff_birthEnv_iff}
    \leanok
    \uses{def:party_sheet}
    For $T\subseteq\Lambda$ and a party $Q$, the sheet $S^{T\to Q}$ hands the sites of $T$ to
    $Q$: its owner is $Q$ on $T$ and $o_S$ elsewhere. For a party $P_\circ$, the set
    $E_S(P_\circ)=\{x:o_S(x)\neq P_\circ\}$ is the set $E$ of
    \cite[Lemma~6.5]{OpenAI2026PolynomialPEPS}, whose hole part is empty on a sheet. Let
    $E=E_S(P_\circ)$ and $U=\Lambda\setminus(T\cup E)$.

    For every $T$ and $Q$, the sites owned by $Q$ in $S^{T\to Q}$ are those of $T$ and those
    owned by $Q$ in $S$; if $Q\neq P_\circ$ the sites owned by $P_\circ$ in $S^{T\to Q}$ are
    those of $U$; and the sites owned by $P_\circ$ in $S^{T\to P_\circ}$ are those of
    $T\cup U$. If moreover every site of $T$ is owned by $P_\circ$ in $S$, then $T$ and $E$ are
    disjoint, $S^{T\to P_\circ}=S$, the sites owned by $P_\circ$ in $S$ are those of $T\cup U$,
    a site lies in $(T\cup E)\setminus T$ exactly when it is not owned by $P_\circ$, every
    party $r\notin\{P_\circ,Q\}$ owns the same sites in $S$ and in $S^{T\to Q}$, the sites of
    $U$ are owned by $P_\circ$ in both sheets, the sites of $T$ by $Q$ in $S^{T\to Q}$, and the
    owners on $E$ are unchanged.
\end{definition}

\begin{definition}[Bases of a party read on regions]
    \label{def:party_sheet_region_coords}
    \lean{QuantumCircuit.Party.Sheet.regionEquiv,
        QuantumCircuit.Party.Sheet.regionEquiv_apply,
        QuantumCircuit.Party.Sheet.regionEquiv_symm_apply,
        QuantumCircuit.Party.Sheet.splitEquiv,
        QuantumCircuit.Party.Sheet.splitEquiv_apply_fst,
        QuantumCircuit.Party.Sheet.splitEquiv_apply_snd,
        QuantumCircuit.Party.Sheet.complSplitEquiv,
        QuantumCircuit.Party.Sheet.system_eq_of_owner_iff,
        QuantumCircuit.Party.Sheet.cast_system_eq_regionEquiv,
        QuantumCircuit.Party.Sheet.birthEquivP, QuantumCircuit.Party.Sheet.birthEquivP',
        QuantumCircuit.Party.Sheet.birthEquivQ, QuantumCircuit.Party.Sheet.birthEquivSelf,
        QuantumCircuit.Party.Sheet.cast_changeOwner_system_iff}
    \leanok
    \uses{def:party_sheet, def:party_change_owner}
    If the sites owned by $p$ in $S$ form the region $R$, restriction identifies $X^S_p$ with
    $[q]^R$; if they form the disjoint union $R_1\sqcup R_2$, it identifies $X^S_p$ with
    $[q]^{R_1}\times[q]^{R_2}$; if every site outside a region $B$ is owned by $P_\circ$, it
    identifies $X^S_{P_\circ}$ with $[q]^{B\cap o_S^{-1}(P_\circ)}\times[q]^{\Lambda\setminus B}$.
    These are relabellings of sites \cite[Section~6]{OpenAI2026PolynomialPEPS}. If $p$ owns the
    same sites in two sheets $S'$ and $S$, then $X^{S'}_p=X^S_p$, and the identification of these
    equal bases is the restriction to the sites of $p$ in $S'$ followed by the inverse of the
    restriction in $S$, site by site.

    In the setting of a birth, with every site of $T$ owned by $P_\circ$, $E=E_S(P_\circ)$ and
    $U=\Lambda\setminus(T\cup E)$, these give the bijections
    \begin{align}
        \beta_P     & : X^S_{P_\circ}\to[q]^T\times[q]^U,                    \notag \\
        \beta_\circ & : X^{S^{T\to P_\circ}}_{P_\circ}\to[q]^T\times[q]^U, \notag   \\
        \beta'_P    & : X^{S^{T\to Q}}_{P_\circ}\to[q]^U,                  \notag   \\
        \beta_Q     & : X^{S^{T\to Q}}_{Q}\to[q]^T\times X^S_Q,
        \label{eq:party_birth_bases}
    \end{align}
    the last two for $Q\neq P_\circ$, where $\beta_Q$ splits a configuration of the sites of $Q$
    after the birth into its part on $T$ and its part on the sites $Q$ owned before. For a party
    $r\notin\{P_\circ,Q\}$, a configuration $y$ of $X^{S^{T\to Q}}$ and a configuration $x$ of
    $X^S$ have the same component at $r$, the two equal bases being identified, exactly when
    $\gamma_{S^{T\to Q}}(y)$ and $\gamma_S(x)$ agree on the sites owned by $r$.
\end{definition}

\subsection{Splitting data and the canonical map of a birth}

\begin{definition}[Splitting data and the canonical map of a birth]
    \label{def:party_splitting_data}
    \lean{ManyBody.SplittingData, ManyBody.SplittingData.error, ManyBody.birthKernel,
        ManyBody.birthEffect, ManyBody.birthOp, ManyBody.configBirthOp}
    \leanok
    \uses{def:sites_config, thm:three_way_cut_coordinates,
        def:canonical_product_purification, def:construction_idle_gates}
    Let $X$ be a system on a finite set of sites $V$, let $A\subseteq B$ be regions and
    $U=V\setminus B$. Write $\mathcal B_A=\C^{X_A}$ and
    $\mathcal B_M=\C^{X_{B\setminus A}\sqcup X_U}$. Splitting data $\sigma$ for $(A,B)$
    \cite[Lemma~6.4]{OpenAI2026PolynomialPEPS} are an isometry
    $V_\sigma:\C^{X_U}\to\mathcal B_A\otimes\mathcal B_M$, a unit vector
    $s\in\C^{X_A}\otimes\mathcal B_A$ and a unit vector
    $s'\in\C^{X_{B\setminus A}}\otimes\mathcal B_M$. Their error on $\Omega\in\C^{X_V}$ is
    \begin{align}
        \varepsilon_\sigma(\Omega)
        =\lVert(1_{X_A\times X_{B\setminus A}}\otimes V_\sigma)\Omega-s\otimes s'\rVert,
        \label{eq:party_splitting_error}
    \end{align}
    with $\Omega$ read in the three-way cut coordinates $((x_A,x_{B\setminus A}),x_U)$ of
    Theorem~\ref{thm:three_way_cut_coordinates} and
    $(s\otimes s')((a,m),(b,b'))=s(a,b)\,s'(m,b')$.

    The kernel of $\sigma$ is the matrix $K$ from $X_A\times X_U$ to the basis of
    $\mathcal B_M$ with
    $K_{b,(a,u)}=\sum_{\beta\in X_A}\overline{s(a,\beta)}\,(V_\sigma)_{(\beta,b),u}$, that is
    $K=\langle s|(1_{X_A}\otimes V_\sigma)$. With $\mathcal V=1\otimes V_\sigma$ on
    $(X_A\times X_{B\setminus A})\times X_U$ and $S_s$ the preparation matrix of $s$ on
    $X_A\times\mathcal B_A$, the effect of a birth is $S_s^\dagger\mathcal V$ and the canonical
    map of the birth is
    \begin{align}
        \mathbf B_\sigma=(S_s^\dagger\mathcal V)^\dagger(S_s^\dagger\mathcal V)
        =\mathcal V^\dagger\bigl(|s\rangle\langle s|\otimes1\bigr)\mathcal V,
        \label{eq:party_birth_map}
    \end{align}
    read on the configurations of $X_V$ along the three-way cut coordinates
    \cite[Lemma~6.5]{OpenAI2026PolynomialPEPS}. In the coordinates $(x_{B\setminus A},(x_A,x_U))$
    it is $\mathbf B_\sigma=1_{X_{B\setminus A}}\otimes K^\dagger K$
    \cite[proof of Lemma~6.5]{OpenAI2026PolynomialPEPS}.
\end{definition}

For a sheet $S$ and a set $T$ of sites owned by $P_\circ$, splitting data for a birth are
splitting data $\sigma$ on $[q]^\Lambda$ for $(T,T\cup E)$ with $E=E_S(P_\circ)$; then
$\mathcal B_T=\C^{[q]^T}$, $\mathcal B_E=\C^{[q]^E\sqcup[q]^U}$ and $(T\cup E)\setminus T=E$.

\begin{theorem}[The canonical map of a birth in the coordinates of the parties]
    \label{thm:party_config_birth_op}
    \lean{QuantumCircuit.Party.Sheet.configBirthOp_apply_changeOwner,
        QuantumCircuit.Party.Sheet.configBirthOp_apply_changeOwner_self}
    \leanok
    \uses{def:party_splitting_data, def:party_sheet_region_coords, def:party_change_owner}
    Let every site of $T$ be owned by $P_\circ$ in the sheet $S$, let $\sigma$ be splitting data
    for the birth with kernel $K$, and use the bijections~(\ref{eq:party_birth_bases}). For
    $Q\neq P_\circ$, a configuration $y$ of $X^{S^{T\to Q}}$ and a configuration $x$ of $X^S$,
    \begin{align}
        \mathbf B_\sigma(\gamma_{S^{T\to Q}}(y),\gamma_S(x))
        =\delta\,(K^\dagger K)\bigl((\beta_Q(y_Q)_1,\beta'_P(y_{P_\circ})),
        \beta_P(x_{P_\circ})\bigr),
        \label{eq:party_birth_entry}
    \end{align}
    where $\delta=1$ if $\beta_Q(y_Q)_2=x_Q$ and $\gamma_{S^{T\to Q}}(y)$, $\gamma_S(x)$ agree on
    the sites of every party $r\notin\{P_\circ,Q\}$, and $\delta=0$ otherwise. For a configuration
    $y$ of $X^{S^{T\to P_\circ}}$,
    $\mathbf B_\sigma(\gamma_{S^{T\to P_\circ}}(y),\gamma_S(x))
        =\delta_\circ\,(K^\dagger K)(\beta_\circ(y_{P_\circ}),\beta_P(x_{P_\circ}))$, where
    $\delta_\circ=1$ if the two configurations agree on the sites of every party $r\neq P_\circ$
    and $\delta_\circ=0$ otherwise.
\end{theorem}
\begin{proof}\leanok
    \uses{def:party_splitting_data, def:party_sheet_region_coords, def:party_change_owner}
    By Definition~\ref{def:party_splitting_data},
    $\mathbf B_\sigma=1_{[q]^E}\otimes K^\dagger K$ in the coordinates $(x_E,(x_T,x_U))$. The
    factor $1_{[q]^E}$ compares the parts of the two configurations on $E$, the sites not owned
    by $P_\circ$: for $Q$ these are the sites it owned before, whose part in $y$ is
    $\beta_Q(y_Q)_2$, and for $r\notin\{P_\circ,Q\}$ the sites $r$ owns in both sheets. The parts
    on $T$ and $U$ are read at $Q$ and $P_\circ$ after the birth and at $P_\circ$ before it,
    through $\beta_Q$, $\beta'_P$ and $\beta_P$. For $Q=P_\circ$ the parts on $T$ and $U$ after
    the birth are both read at $P_\circ$ through $\beta_\circ$.
\end{proof}

\subsection{Births and deaths}

In the rest of this subsection $S$ is a sheet, every site of $T$ is owned by $P_\circ$,
$E=E_S(P_\circ)$, $U=\Lambda\setminus(T\cup E)$, and $\sigma$ are splitting data for the birth
with isometry $V_\sigma:\C^{[q]^U}\to\mathcal B_T\otimes\mathcal B_E$, unit vector
$s\in\C^{[q]^T}\otimes\mathcal B_T$ and kernel $K$.

\begin{definition}[Birth words and death chains]
    \label{def:party_birth_word}
    \lean{QuantumCircuit.Party.Sheet.birthPartyMatrix, QuantumCircuit.Party.Sheet.birthFamily,
        QuantumCircuit.Party.Sheet.birthWord, QuantumCircuit.Party.Sheet.deathFamily,
        QuantumCircuit.Party.Sheet.deathChain}
    \leanok
    \uses{def:party_splitting_data, def:party_sheet_region_coords, def:party_word,
        def:party_add_pair, def:party_pair_gates, def:party_circuit}
    Let $Q\neq P_\circ$. The contraction of the birth at $P_\circ$ is
    $M_\sigma=V_\sigma^\dagger(1_{\mathcal B_T}\otimes K)$, from
    $\mathcal B_T\otimes\C^{[q]^T}\otimes\C^{[q]^U}$ to $\C^{[q]^U}$.

    The birth word from $X^S$ to $X^{S^{T\to Q}}$ is the normalized pair source $s$ on a fresh
    copy of $[q]^T$ added at $Q$ and a fresh factor $\mathcal B_T$ added at $P_\circ$, followed
    by one local operation at every party: at $Q$ the matrix of the bijection $\beta_Q^{-1}$,
    so that the new copy of $T$ joins the sites of $Q$; at $P_\circ$ the matrix $M_\sigma$, read
    on the bases through $\beta_P$ and $\beta'_P$, so that the copy of $T$ held by $P_\circ$
    before the birth leaves $P_\circ$; and the identification of equal bases at every other
    party. In \cite[proof of Lemma~6.5]{OpenAI2026PolynomialPEPS} the contraction $K$ is applied
    before the pair source; it acts on factors other than those of the source, and here it is
    applied after the source, together with $V_\sigma^\dagger$, in one local operation at
    $P_\circ$.

    The death chain from $X^{S^{T\to Q}}$ to $X^S$ is the monomial made of one local operation
    at every party, the matrix of $\beta_Q$ at $Q$, the adjoint
    $M_\sigma^\dagger=(1_{\mathcal B_T}\otimes K^\dagger)V_\sigma$ at $P_\circ$ read through
    the same bijections, and the identification of equal bases elsewhere, followed by the
    normalized pair effect $\langle s|$ on the copy of $T$ at $Q$ and on $\mathcal B_T$ at
    $P_\circ$. Since $M_\sigma^\dagger$ acts on factors other than those of the pair effect, it is
    applied before the effect.
\end{definition}

\begin{theorem}[The contraction of a birth]
    \label{thm:party_birth_contraction}
    \lean{QuantumCircuit.Party.Sheet.norm_birthPartyMatrix_le_one,
        QuantumCircuit.Party.Sheet.sum_birthPartyMatrix_mul}
    \leanok
    \uses{def:party_birth_word}
    Let $T$, $U$, $B_T$ and $B_E$ be finite sets, $V$ a matrix from $\C^U$ to
    $\C^{B_T}\otimes\C^{B_E}$, $s$ a vector on $T\times B_T$, $K=\langle s|(1_T\otimes V)$ and
    $M=V^\dagger(1_{B_T}\otimes K)$. Then for all $t'\in T$, $u\in U$ and $(t,u')\in T\times U$,
    \begin{align}
        \sum_{b\in B_T}M_{u,(b,(t,u'))}\,s(t',b)=(K^\dagger K)_{(t',u),(t,u')},
        \label{eq:party_birth_contraction}
    \end{align}
    and if $V$ is an isometry and $s$ a unit vector, then $\lVert M\rVert\leq1$.
\end{theorem}
\begin{proof}\leanok
    \uses{def:party_splitting_data}
    Expanding the identity factor,
    $M_{u,(b,\tau)}=\sum_{b'}\overline{V_{(b,b'),u}}\,K_{b',\tau}$, and
    $\sum_b\overline{V_{(b,b'),u}}\,s(t',b)=\overline{K_{b',(t',u)}}$, which
    gives~(\ref{eq:party_birth_contraction}) after exchanging the two sums. For the norm,
    $\lVert V^\dagger\rVert\leq1$, and $\lVert1\otimes K\rVert\leq\lVert K\rVert\leq1$ because $K$
    is the effect $S_s^\dagger\mathcal V$ of Definition~\ref{def:party_splitting_data} with a
    one-point middle region, a product of two contractions.
\end{proof}

\begin{theorem}[Birth words and death chains implement the canonical map]
    \label{thm:party_birth_word}
    \lean{QuantumCircuit.Party.Sheet.op_birthWord,
        QuantumCircuit.Party.Sheet.isAllowed_birthWord,
        QuantumCircuit.Party.Sheet.usesOnly_birthWord,
        QuantumCircuit.Party.Sheet.sourceCount_birthWord,
        QuantumCircuit.Party.Sheet.op_deathChain,
        QuantumCircuit.Party.Sheet.isAllowed_deathChain,
        QuantumCircuit.Party.Sheet.usesOnly_deathChain,
        QuantumCircuit.Party.Sheet.forallGate_isRight_deathChain,
        QuantumCircuit.Party.Sheet.numGates_deathChain}
    \leanok
    \uses{def:party_birth_word, def:party_allowed, def:party_word, def:party_gate_counts}
    Let $Q\neq P_\circ$ and $S'=S^{T\to Q}$. The birth word is allowed, uses only $P_\circ$
    and $Q$, has exactly one pair source, and its operator is the canonical map in canonical
    coordinates: $\operatorname{op}(y,x)=\mathbf B_\sigma(\gamma_{S'}(y),\gamma_S(x))$ for
    configurations $y$ of $X^{S'}$ and $x$ of $X^S$. The death chain is allowed, uses only
    $P_\circ$ and $Q$, has exactly one gate, which is a pair effect, and its operator is
    $\operatorname{op}(x,y)=\mathbf B_\sigma(\gamma_S(x),\gamma_{S'}(y))$.
\end{theorem}
\begin{proof}\leanok
    \uses{thm:party_birth_contraction, thm:party_config_birth_op, def:party_birth_word,
        def:party_splitting_data}
    The operator of the birth word is the product over the parties of the local operations
    applied after the preparation of $s$. At the configurations $y$, $x$ and the fresh pair
    $(a,b)$, the factor at $Q$ is $1$ exactly when $\beta_Q(y_Q)=(a,x_Q)$, the factor at $P_\circ$
    is $M_\sigma(\beta'_P(y_{P_\circ}),(b,\beta_P(x_{P_\circ})))$, and the factor at every other
    party is $1$ exactly when the configurations agree on its sites. The sum over $a$ leaves
    $a=\beta_Q(y_Q)_1$, and the sum over $b$ against $s(a,b)$ gives $K^\dagger K$ by
    Theorem~\ref{thm:party_birth_contraction}; the result is the right side
    of~(\ref{eq:party_birth_entry}) in Theorem~\ref{thm:party_config_birth_op}. The death chain
    gives the same computation with the adjoint factors and the effect $\langle s|$, hence the
    complex conjugate of $\mathbf B_\sigma(\gamma_{S'}(y),\gamma_S(x))$, which is
    $\mathbf B_\sigma(\gamma_S(x),\gamma_{S'}(y))$ because $\mathbf B_\sigma$ is Hermitian
    by~(\ref{eq:party_birth_map}). Both are allowed: $s$ is a unit vector, the matrices of
    bijections are contractions and $\lVert M_\sigma\rVert\leq1$ by
    Theorem~\ref{thm:party_birth_contraction}. Outside $\{P_\circ,Q\}$ the local operations are
    identifications of equal bases, and the counts hold by definition.
\end{proof}

\begin{definition}[Births and deaths at one party]
    \label{def:party_private_birth}
    \lean{QuantumCircuit.Party.Sheet.privateBirthFamily,
        QuantumCircuit.Party.Sheet.privateBirthWord,
        QuantumCircuit.Party.Sheet.privateDeathChain}
    \leanok
    \uses{def:party_splitting_data, def:party_sheet_region_coords, def:party_word,
        def:party_pair_gates, def:party_circuit}
    When the party receiving $T$ is $P_\circ$ itself, no pair source is needed
    \cite[proof of Lemma~6.5]{OpenAI2026PolynomialPEPS}. The private birth word from $X^S$ to
    $X^{S^{T\to P_\circ}}$ is the word without gates made of one local operation at every party:
    the matrix $K^\dagger K$ at $P_\circ$, read on the bases through $\beta_P$ and $\beta_\circ$,
    and the identification of equal bases at every other party. The private death chain from
    $X^{S^{T\to P_\circ}}$ to $X^S$ is the monomial without gates made of the adjoints of these
    local operations.
\end{definition}

\begin{theorem}[Births and deaths at one party implement the canonical map]
    \label{thm:party_private_birth}
    \lean{QuantumCircuit.Party.Sheet.op_privateBirthWord,
        QuantumCircuit.Party.Sheet.isAllowed_privateBirthWord,
        QuantumCircuit.Party.Sheet.usesOnly_privateBirthWord,
        QuantumCircuit.Party.Sheet.sourceCount_privateBirthWord,
        QuantumCircuit.Party.Sheet.op_privateDeathChain,
        QuantumCircuit.Party.Sheet.isAllowed_privateDeathChain,
        QuantumCircuit.Party.Sheet.usesOnly_privateDeathChain,
        QuantumCircuit.Party.Sheet.numGates_privateDeathChain}
    \leanok
    \uses{def:party_private_birth, def:party_allowed, def:party_word, def:party_gate_counts}
    Let $S'=S^{T\to P_\circ}$. The private birth word is allowed, uses only $P_\circ$, has no
    pair source, and its operator is
    $\operatorname{op}(y,x)=\mathbf B_\sigma(\gamma_{S'}(y),\gamma_S(x))$. The private death chain
    is allowed, uses only $P_\circ$, has no gate, and its operator is
    $\operatorname{op}(x,y)=\mathbf B_\sigma(\gamma_S(x),\gamma_{S'}(y))$.
\end{theorem}
\begin{proof}\leanok
    \uses{thm:party_config_birth_op, def:party_private_birth, def:party_splitting_data}
    The operator of the private birth word is the product over the parties of its local
    operations: $(K^\dagger K)(\beta_\circ(y_{P_\circ}),\beta_P(x_{P_\circ}))$ at $P_\circ$ times
    the agreement of the configurations on the sites of every other party, which is the second
    formula of Theorem~\ref{thm:party_config_birth_op}. The operator of the death chain is the
    adjoint of this product, and $\mathbf B_\sigma$ is Hermitian by~(\ref{eq:party_birth_map}).
    Both are allowed because $\lVert K\rVert\leq1$ (see the proof of
    Theorem~\ref{thm:party_birth_contraction}), so $\lVert K^\dagger K\rVert\leq1$, and the
    identifications of equal bases are contractions; the remaining statements hold by
    definition.
\end{proof}

\begin{theorem}[Homogeneous birth and death for given splitting data]
    \label{thm:party_birth}
    \lean{QuantumCircuit.Party.Sheet.exists_birthWord_of_environment,
        QuantumCircuit.Party.Sheet.exists_deathChain_of_environment,
        QuantumCircuit.Party.Sheet.exists_birthWord_of_splittingData,
        QuantumCircuit.Party.Sheet.exists_deathChain_of_splittingData}
    \leanok
    \uses{def:party_change_owner, def:party_splitting_data, def:party_word,
        def:party_pair_gates, def:party_allowed, def:party_gate_counts}
    Let $S$ be a sheet, $P_\circ$ and $Q$ parties, not necessarily distinct, $T\subseteq\Lambda$
    a set of sites owned by $P_\circ$, and $E\subseteq\Lambda$ a set disjoint from $T$ such that
    every site outside $T\cup E$ is owned by $P_\circ$; sites of $E$ may be owned by $P_\circ$.
    Let $\sigma$ be splitting data on $[q]^\Lambda$ for $(T,T\cup E)$, let
    $\Omega\in\C^{[q]^\Lambda}$ and $S'=S^{T\to Q}$.
    \begin{enumerate}
        \item There is an allowed word $w$ from $X^S$ to $X^{S'}$ that uses only $P_\circ$ and $Q$,
              has at most one pair source, whose operator is
              $\operatorname{op}(w)(y,x)=\mathbf B_\sigma(\gamma_{S'}(y),\gamma_S(x))$, and with
              $\lVert\operatorname{op}(w)\Omega_S-\Omega_{S'}\rVert\leq\varepsilon_\sigma(\Omega)$.
        \item There is an allowed monomial $M$ from $X^{S'}$ to $X^S$ that uses only $P_\circ$
              and $Q$, whose gates are pair effects and at most one, whose operator is
              $\operatorname{op}(M)(x,y)=\mathbf B_\sigma(\gamma_S(x),\gamma_{S'}(y))$, and with
              $\lVert\operatorname{op}(M)\Omega_{S'}-\Omega_S\rVert\leq\varepsilon_\sigma(\Omega)$.
    \end{enumerate}
    In particular both hold for $E=E_S(P_\circ)$. This is the part of
    \cite[Lemma~6.5]{OpenAI2026PolynomialPEPS} that holds for any splitting data: the birth and
    the death involve only $P_\circ$ and $Q$, use at most one normalized pair source or pair
    effect, and move the reference vector by at most the error of the splitting data. The
    hypothesis of the lemma on the mutual information of $T$ and $E$ enters only through the
    existence of splitting data of small error \cite[Lemma~6.4]{OpenAI2026PolynomialPEPS},
    which is not part of this statement.
\end{theorem}
\begin{proof}\leanok
    \uses{def:party_change_owner, def:party_splitting_data, def:party_sheet}
    The canonical coordinates are bijections, so a word or monomial whose operator is
    $\mathbf B_\sigma$ in canonical coordinates moves $\Omega_S$ to within
    $\lVert\mathbf B_\sigma\Omega-\Omega\rVert$ of $\Omega_{S'}$. With the projection
    $\Pi=S_sS_s^\dagger$, which fixes $s\otimes s'$,
    $(\mathbf B_\sigma-1)\Omega=\mathcal V^\dagger(\Pi-1)(\mathcal V\Omega-s\otimes s')$ because
    $\mathcal V$ is an isometry, and $\lVert\mathcal V^\dagger\rVert\leq1$,
    $\lVert\Pi-1\rVert\leq1$ give
    $\lVert\mathbf B_\sigma\Omega-\Omega\rVert\leq\varepsilon_\sigma(\Omega)$
    by~(\ref{eq:party_splitting_error}). For $Q=P_\circ$ the sheet is unchanged,
    $S^{T\to P_\circ}=S$, and the private birth word and private death chain of
    Definition~\ref{def:party_private_birth} serve, modified so that the local operation at
    $P_\circ$ acts as the identity on the sites of $E$ owned by $P_\circ$. For $Q\neq P_\circ$ the
    birth word and death chain of Definition~\ref{def:party_birth_word} serve, with the same
    modification of the local operation at $P_\circ$. The computation of
    Theorems~\ref{thm:party_birth_word} and~\ref{thm:party_private_birth}, with the sites of
    $E$ owned by $P_\circ$ compared by the identity at $P_\circ$, shows that their operators are
    $\mathbf B_\sigma$ in canonical coordinates, and the counts and the norms of the local
    operations are as there. For $E=E_S(P_\circ)$, the set $T$ is disjoint from $E$ and the sites outside
    $T\cup E$ are owned by $P_\circ$ by Definition~\ref{def:party_change_owner}.
\end{proof}

\subsection{Two-sheet exchange}

\begin{definition}[Exchange of two sheets inside a region]
    \label{def:party_exchange_sets}
    \lean{QuantumCircuit.Party.Sheet.swapIn, QuantumCircuit.Party.Sheet.exchangeEnv,
        QuantumCircuit.Party.Sheet.mem_exchangeEnv,
        QuantumCircuit.Party.Sheet.owner_eq_of_notMem_exchangeEnv,
        QuantumCircuit.Party.Sheet.sdiff_subset_compl_inter_union_sdiff,
        QuantumCircuit.Party.Sheet.swapIn_owner_eq_of_notMem,
        QuantumCircuit.Party.Sheet.swapIn_owner_eq_of_notMem_union,
        QuantumCircuit.Party.Sheet.swapIn_owner_eq_of_notMem_union'}
    \leanok
    \uses{def:party_sheet}
    For sheets $S_1,S_2$ and a region $Y$, the sheet $S_1\ast_YS_2$ has the owner $o_{S_2}$ on
    $Y$ and $o_{S_1}$ outside $Y$. The exchange inside $Y$ replaces $(S_1,S_2)$ by
    $(S_1',S_2')=(S_1\ast_YS_2,S_2\ast_YS_1)$ \cite[Lemma~6.6]{OpenAI2026PolynomialPEPS}. For a
    party $P_\circ$, the set $Z_\circ$ of sites whose two owners are not both $P_\circ$ is the set
    $Z$ of the lemma without its hole parts, which are empty on sheets; so let $Z$ be any set of
    sites containing $Z_\circ$, that is, outside of which both owners are $P_\circ$. Put
    $T=Z\cap Y$, $E=Z\setminus Y$, $A=Y\setminus Z$ and $U=\Lambda\setminus Z$. Then
    $T\cup E=Z$, $A\subseteq\Lambda\setminus(T\cup E)$, and every site outside $T\cup E$ is owned
    by $P_\circ$ on both sheets after the exchange.
\end{definition}

\begin{definition}[The map of an exchange]
    \label{def:party_exchange_map}
    \lean{QuantumCircuit.Party.Sheet.exchangeOp}
    \leanok
    \uses{def:party_exchange_sets, def:party_splitting_data, def:sites_pair_swap}
    Let $\sigma$ be splitting data on $[q]^\Lambda$ for $(T,T\cup E)$, with isometry
    $V_\sigma:\C^{[q]^U}\to\mathcal B_T\otimes\mathcal B_E$. On two copies
    $\C^{[q]^\Lambda\times[q]^\Lambda}$, let $F_R$ be the exchange of the two copies on a region
    $R$, and let the buffer correction $D_U$ act, in the three-way cut coordinates of both
    copies, as the identity on the two copies of $[q]^T\times[q]^E$ and as
    $(V_\sigma\otimes V_\sigma)^\dagger F_{\mathcal B_T}(V_\sigma\otimes V_\sigma)$ on the two
    copies of $[q]^U$, where $F_{\mathcal B_T}$ exchanges the $\mathcal B_T$ factors of
    $(\mathcal B_T\otimes\mathcal B_E)^{\otimes2}$. The map of the exchange is
    $C_\sigma=D_UF_AF_Y$ \cite[proof of Lemma~6.6]{OpenAI2026PolynomialPEPS}.
\end{definition}

\begin{definition}[Renaming and corrections of an exchange]
    \label{def:party_exchange_words}
    \lean{QuantumCircuit.Party.Sheet.swapInEquiv, QuantumCircuit.Party.Sheet.renameWord,
        QuantumCircuit.Party.Sheet.exchangeEquivP, QuantumCircuit.Party.Sheet.correctionAt,
        QuantumCircuit.Party.Sheet.correctionWord}
    \leanok
    \uses{def:party_exchange_map, def:party_sheet_region_coords, def:party_relabel,
        def:party_word, def:party_circuit}
    At every party $p$, the bijection $\rho_p$ from $X^{S_1'}_p\times X^{S_2'}_p$ to
    $X^{S_1}_p\times X^{S_2}_p$ moves the coordinate of a site of $Y$ to the other sheet and keeps
    the coordinates of the other sites. The renaming $\mathcal R$ is the word relabelling the
    bases along $\rho$; every register keeps its party. Restriction identifies
    $X^{S_1'}_{P_\circ}\times X^{S_2'}_{P_\circ}$ with the product of the configurations of the
    sites of $Z$ owned by $P_\circ$ on each sheet and of two copies of $[q]^U$. The correction at
    $P_\circ$ is the identity on the first factor tensored with
    \begin{align}
        W_\sigma=(V_\sigma\otimes V_\sigma)^\dagger F_{\mathcal B_T}(V_\sigma\otimes V_\sigma)F_A
        \label{eq:party_correction}
    \end{align}
    on the two copies of $[q]^U$, with $F_A$ exchanging the parts on $A\subseteq U$ of the two
    copies. The correction word is the word without gates whose local operation is the
    correction at $P_\circ$ and the identity at every other party.
\end{definition}

\begin{theorem}[Renaming and corrections in canonical coordinates]
    \label{thm:party_exchange_words}
    \lean{QuantumCircuit.Party.Sheet.twoConfigEquiv_symm_permCongr_pairSwap,
        QuantumCircuit.Party.Sheet.op_renameWord,
        QuantumCircuit.Party.Sheet.isAllowed_renameWord,
        QuantumCircuit.Party.Sheet.sourceCount_renameWord,
        QuantumCircuit.Party.Sheet.norm_correctionAt_le_one,
        QuantumCircuit.Party.Sheet.op_correctionWord,
        QuantumCircuit.Party.Sheet.isAllowed_correctionWord,
        QuantumCircuit.Party.Sheet.usesOnly_correctionWord,
        QuantumCircuit.Party.Sheet.sourceCount_correctionWord,
        QuantumCircuit.Party.Sheet.op_comp_renameWord_correctionWord}
    \leanok
    \uses{def:party_exchange_words, def:party_allowed}
    In the setting of Definition~\ref{def:party_exchange_words}:
    \begin{enumerate}
        \item The exchange of the two copies on $Y$, read from the canonical coordinates of
              $(S_1',S_2')$ to those of $(S_1,S_2)$, is the bijection $(\rho_p)_p$, that is
              $\gamma_{S_1,S_2}^{-1}\circ F_Y\circ\gamma_{S_1',S_2'}=(\rho_p)_p$. Hence the operator
              of $\mathcal R$ is $F_Y$ in canonical coordinates, and $\mathcal R$ is allowed and
              has no pair source.
        \item The correction at $P_\circ$ is a contraction. The correction word is allowed, uses
              only $P_\circ$, has no pair source, and its operator in the canonical coordinates of
              $(S_1',S_2')$ is $D_UF_A$.
        \item The renaming followed by the correction word has the operator $C_\sigma$, read from
              the canonical coordinates of $(S_1,S_2)$ to those of $(S_1',S_2')$.
    \end{enumerate}
\end{theorem}
\begin{proof}\leanok
    \uses{def:party_exchange_words, def:party_exchange_sets, def:party_relabel,
        def:sites_pair_swap}
    The first statement is checked site by site: a site of $Y$ in the first sheet after the
    exchange is read at its owner in the second sheet before it, and conversely, while the other
    sites stay on their sheets. The operator of a relabelling is the matrix of the bijection of
    configurations, which gives $F_Y$. The norm of the correction is at most
    $\lVert W_\sigma\rVert\leq1$, since $V_\sigma\otimes V_\sigma$ is an isometry and the
    exchanges are permutation matrices. Since $A\subseteq U$, the map $D_UF_A$ acts as the
    identity on the two copies of the configurations of $Z$ and as $W_\sigma$ on the two copies
    of $[q]^U$. Every site outside $Z$ is owned by $P_\circ$ on both sheets after the exchange,
    so the parts of two configurations on the sites of the other parties are compared by the
    identity local operations, and the parts on the sites of $Z$ owned by $P_\circ$ by the
    identity factor of the correction; this gives the operator of the correction word. The last
    statement is the product of the first two.
\end{proof}

\begin{theorem}[Two-sheet exchange for given splitting data]
    \label{thm:party_exchange}
    \lean{QuantumCircuit.Party.Sheet.exists_exchangeWord_of_splittingData,
        QuantumCircuit.Party.Sheet.exchangeOp_eq}
    \leanok
    \uses{def:party_exchange_map, def:party_exchange_words, def:party_allowed, def:party_word}
    Let $S_1,S_2$ be sheets, $Y$ a region, $P_\circ$ a party, $Z$ a set of sites outside of
    which both owners are $P_\circ$, with $T$, $E$, $A$, $U$ as in
    Definition~\ref{def:party_exchange_sets}, and let $\sigma$ be splitting data on
    $[q]^\Lambda$ for $(T,T\cup E)$. Then
    \begin{align}
        C_\sigma=D_UF_AF_Y=D_UF_T,
        \label{eq:party_exchange_net}
    \end{align}
    and there is an allowed word $c$ from $X^{S_1'}\times X^{S_2'}$ to itself that uses only
    $P_\circ$ and has no pair source, such that the renaming $\mathcal R$ followed by $c$ has the
    operator $C_\sigma$, read from the canonical coordinates of $(S_1,S_2)$ to those of
    $(S_1',S_2')$, and, for every unit vector $\Omega\in\C^{[q]^\Lambda}$,
    \begin{align}
        \bigl\lVert\operatorname{op}(\mathcal R;c)(\Omega_{S_1}\otimes\Omega_{S_2})
        -\Omega_{S_1'}\otimes\Omega_{S_2'}\bigr\rVert\leq4\,\varepsilon_\sigma(\Omega).
        \label{eq:party_exchange_error}
    \end{align}
    This is \cite[Lemma~6.6]{OpenAI2026PolynomialPEPS} for given splitting data: the exchange is
    a relabelling of the bases of the parties whose sites change sheet, no register changing
    party, followed by corrections private to $P_\circ$. It has no gate, so it is a bounded
    change for every set of participants (Theorem~\ref{thm:party_resources_bounded}).
\end{theorem}
\begin{proof}\leanok
    \uses{thm:party_exchange_words, def:party_sheet, def:sites_pair_swap}
    Since $A\subseteq Y$, the exchanges compose to $F_AF_Y=F_{Y\setminus A}$, and
    $Y\setminus(Y\setminus Z)=Z\cap Y=T$, which is~(\ref{eq:party_exchange_net}). The word $c$ is
    the correction word, and Theorem~\ref{thm:party_exchange_words} gives its properties and the
    operator of $\mathcal R$ followed by $c$. By Definition~\ref{def:party_sheet},
    $\Omega_{S_1}\otimes\Omega_{S_2}$ and $\Omega_{S_1'}\otimes\Omega_{S_2'}$ are $\Omega\otimes\Omega$
    in the canonical coordinates, so the left side of~(\ref{eq:party_exchange_error}) is
    $\lVert D_UF_T(\Omega\otimes\Omega)-\Omega\otimes\Omega\rVert$. With
    $\mathcal V=1\otimes V_\sigma$ and $\zeta=s\otimes s'$,
    $D_UF_T=(\mathcal V^{\otimes2})^\dagger F_{T\mathcal B_T}\mathcal V^{\otimes2}$, where
    $F_{T\mathcal B_T}$ exchanges the parts on $[q]^T\times\mathcal B_T$ of the two copies, and
    $F_{T\mathcal B_T}$ fixes $\zeta\otimes\zeta$, so
    $D_UF_T(\Omega\otimes\Omega)-\Omega\otimes\Omega=(\mathcal V^{\otimes2})^\dagger
        (F_{T\mathcal B_T}-1)\bigl((\mathcal V\Omega)^{\otimes2}-\zeta\otimes\zeta\bigr)$. Since
    $\lVert F_{T\mathcal B_T}-1\rVert\leq2$ and
    $\lVert(\mathcal V\Omega)^{\otimes2}-\zeta\otimes\zeta\rVert\leq
        2\lVert\mathcal V\Omega-\zeta\rVert$ for unit vectors, the bound
    is~$4\,\varepsilon_\sigma(\Omega)$ \cite[proof of Lemma~6.6]{OpenAI2026PolynomialPEPS}.
\end{proof}

\subsection{Stacks of sheets}

In this subsection the sheets form an arbitrary set $\kappa$; every $F\in\kappa$ carries a system
$X^F$ on the parties and, where needed, a reference vector $r_F\in\mathcal H_{X^F}$. Ownership
sheets with $X^S$ and $r_S=\Omega_S$ are one instance.

\begin{definition}[Stacks of sheets and lifted changes]
    \label{def:party_sheet_stack}
    \lean{QuantumCircuit.Party.stackSystem, QuantumCircuit.Party.stackVec,
        QuantumCircuit.Party.norm_stackVec, QuantumCircuit.Party.norm_stackVec_le,
        QuantumCircuit.Party.framesVec, QuantumCircuit.Party.norm_framesVec,
        QuantumCircuit.Party.liftOne, QuantumCircuit.Party.liftTwo,
        QuantumCircuit.Party.nonprivateCount_liftOne,
        QuantumCircuit.Party.forall_participants_liftOne,
        QuantumCircuit.Party.isExpansionBounded_liftOne,
        QuantumCircuit.Party.isMonomialBounded_liftOne,
        QuantumCircuit.Party.isSourceBounded_liftOne,
        QuantumCircuit.Party.nonprivateCount_liftTwo,
        QuantumCircuit.Party.forall_participants_liftTwo,
        QuantumCircuit.Party.isExpansionBounded_liftTwo,
        QuantumCircuit.Party.isMonomialBounded_liftTwo,
        QuantumCircuit.Party.isSourceBounded_liftTwo}
    \leanok
    \uses{def:party_original_frame_sequence}
    For a finite list $\mathbf F=(F_1,\dots,F_n)$ of sheets and an idle system $W$, the stack
    system $\mathrm{St}(\mathbf F;W)$ has at the party $p$ the basis
    $X^{F_1}_p\times(X^{F_2}_p\times\cdots\times(X^{F_n}_p\times W_p))$, and $W_p$ for the empty
    list. For $\xi\in\mathcal H_W$ the reference vector of the stack is
    $\omega_{\mathbf F}(\xi)=r_{F_1}\otimes(r_{F_2}\otimes\cdots\otimes(r_{F_n}\otimes\xi))$
    \cite[Proposition~7.1]{OpenAI2026PolynomialPEPS}. Its norm is
    $\lVert\omega_{\mathbf F}(\xi)\rVert=\prod_i\lVert r_{F_i}\rVert\,\lVert\xi\rVert$, at most
    $\lVert\xi\rVert$ when every $\lVert r_F\rVert\leq1$. For $W$ the system with one basis vector
    at every party and $\xi$ that vector, $\omega_{\mathbf F}=\omega_{\mathbf F}(\xi)$ has norm
    $\prod_i\lVert r_{F_i}\rVert$.

    Write $(\mathbf G,F,\mathbf H)$ for the list $\mathbf G$ followed by $F$ and then $\mathbf H$.
    For an original circuit $c$ from $X^F$ to $X^{F'}$, the lift $L_1(c)$ from
    $\mathrm{St}((\mathbf G,F,\mathbf H);W)$ to $\mathrm{St}((\mathbf G,F',\mathbf H);W)$ is, for
    $\mathbf G$ empty, the relabelling exchanging the two factors of
    $X^F\times\mathrm{St}(\mathbf H;W)$, then $c$ framed by $\mathrm{St}(\mathbf H;W)$, then the
    relabelling exchanging the two factors back; for $\mathbf G=(G,\mathbf G')$ it is the lift
    for $\mathbf G'$ framed by $X^G$. The lift $L_2(c)$ of an original circuit $c$ from
    $X^{F_1}\times X^{F_2}$ to $X^{F_1'}\times X^{F_2'}$, from
    $\mathrm{St}((\mathbf G,F_1,F_2,\mathbf H);W)$ to $\mathrm{St}((\mathbf G,F_1',F_2',\mathbf H);W)$,
    is defined in the same way, with relabellings moving $\mathrm{St}(\mathbf H;W)$ in front of
    the product of the two sheets and back. Lifting keeps the number of nonprivate gates, every
    condition on the participants of the gates, and the expansion, monomial and pair-source
    bounds.
\end{definition}

\begin{theorem}[The reference error of a lifted change]
    \label{thm:party_stack_lift_error}
    \lean{QuantumCircuit.Party.norm_op_liftOne_stackVec_sub_le,
        QuantumCircuit.Party.norm_op_liftTwo_stackVec_sub_le}
    \leanok
    \uses{def:party_sheet_stack}
    Let $\lVert r_G\rVert\leq1$ for every sheet $G$ and $\lVert\xi\rVert\leq1$. For an original
    circuit $c$ from $X^F$ to $X^{F'}$,
    \begin{align}
        \bigl\lVert\operatorname{op}(L_1(c))\,\omega_{(\mathbf G,F,\mathbf H)}(\xi)
        -\omega_{(\mathbf G,F',\mathbf H)}(\xi)\bigr\rVert
        \leq\lVert\operatorname{op}(c)r_F-r_{F'}\rVert,
    \end{align}
    and for an original circuit $c$ from $X^{F_1}\times X^{F_2}$ to $X^{F_1'}\times X^{F_2'}$ the
    lift $L_2(c)$ moves $\omega_{(\mathbf G,F_1,F_2,\mathbf H)}(\xi)$ to within
    $\lVert\operatorname{op}(c)(r_{F_1}\otimes r_{F_2})-r_{F_1'}\otimes r_{F_2'}\rVert$ of
    $\omega_{(\mathbf G,F_1',F_2',\mathbf H)}(\xi)$ \cite[Proposition~7.1]{OpenAI2026PolynomialPEPS}.
\end{theorem}
\begin{proof}\leanok
    \uses{thm:party_frame_prod_vec, def:party_sheet_stack}
    By induction on $\mathbf G$. For $\mathbf G$ empty, write $\zeta=\omega_{\mathbf H}(\xi)$, so
    $\lVert\zeta\rVert\leq\lVert\xi\rVert\leq1$. By Theorem~\ref{thm:party_frame_prod_vec} the first
    relabelling maps $r_F\otimes\zeta$ to $\zeta\otimes r_F$, the framed circuit maps it to
    $\zeta\otimes\operatorname{op}(c)r_F$, and the last relabelling to
    $\operatorname{op}(c)r_F\otimes\zeta$; the difference with $r_{F'}\otimes\zeta$ is
    $(\operatorname{op}(c)r_F-r_{F'})\otimes\zeta$, of norm at most
    $\lVert\operatorname{op}(c)r_F-r_{F'}\rVert$. For two sheets the relabellings map
    $r_{F_1}\otimes(r_{F_2}\otimes\zeta)$ to $\zeta\otimes(r_{F_1}\otimes r_{F_2})$ and back, up to
    the reassociation of the factors, which preserves norms. For $\mathbf G=(G,\mathbf G')$, the
    bound of Theorem~\ref{thm:party_frame_prod_vec} for the circuit framed by $X^G$, with
    $\lVert r_G\rVert\leq1$, reduces the statement to that for $\mathbf G'$.
\end{proof}

\begin{theorem}[Reordering the sheets is a relabelling]
    \label{thm:party_stack_perm}
    \lean{QuantumCircuit.Party.exists_relabel_stackVec_of_perm,
        QuantumCircuit.Party.perm_map_update_cons, QuantumCircuit.Party.perm_map_cons}
    \leanok
    \uses{def:party_sheet_stack}
    If the list $\mathbf F'$ is a permutation of the list $\mathbf F$, there are bijections
    $e_p:\mathrm{St}(\mathbf F';W)_p\to\mathrm{St}(\mathbf F;W)_p$ such that the relabelling
    along $e$ maps $\omega_{\mathbf F}(\xi)$ to $\omega_{\mathbf F'}(\xi)$ for every
    $\xi\in\mathcal H_W$. Moreover, for a map $a$ on a set of indices, an element $x$ of a list
    $l$ of indices, and values $m$ and $f$, the list $(m,a(l))$ is a permutation of
    $(a(x),m,a(l\setminus x))$, where $l\setminus x$ is $l$ with one occurrence of $x$ removed;
    if $l$ has no repetition and $a'$ is $a$ with the value at $x$ replaced by $f$, then
    $(m,a'(l))$ is a permutation of $(f,m,a(l\setminus x))$.
\end{theorem}
\begin{proof}\leanok
    \uses{thm:party_frame_prod_vec}
    By induction on the permutation. The empty list needs the identity. Adding the same sheet
    $F$ in front of both lists keeps the identity on $X^F$ beside the given bijections. Exchanging
    the first two sheets is the reassociation and commutation of the factors
    $X^{F}_p\times(X^{G}_p\times R)\to X^{G}_p\times(X^{F}_p\times R)$, which exchanges $r_F$ and
    $r_G$ in the product vector. A composite of permutations is handled by composing the
    bijections. The list statements follow from the permutation of $l$ to $(x,l\setminus x)$,
    mapped by $a$ or $a'$, followed by the exchange of the first two entries; when $l$ has no
    repetition, $a'$ and $a$ agree on $l\setminus x$.
\end{proof}

\subsection{Bounded changes and the initialization}

\begin{definition}[Bounded changes]
    \label{def:party_bounded_change}
    \lean{QuantumCircuit.Party.OriginalCircuit.IsBoundedChange,
        QuantumCircuit.Party.OriginalCircuit.IsBoundedChange.mono,
        QuantumCircuit.Party.HasBoundedChange, QuantumCircuit.Party.HasSourceBoundedChange,
        QuantumCircuit.Party.HasSourceBoundedChange.hasBoundedChange,
        QuantumCircuit.Party.HasBoundedChange.mono,
        QuantumCircuit.Party.HasSourceBoundedChange.mono}
    \leanok
    \uses{def:party_original_frame_sequence, def:party_sheet_stack}
    Let $c$ be an original circuit from $X$ to $Y$, $S$ a set of parties, $r,K\in\N$,
    $\varepsilon\in\R$, $x\in\mathcal H_X$ and $y\in\mathcal H_Y$. The circuit $c$ is a bounded
    change from $x$ to $y$ with participants in $S$, bounds $r$, $K$ and error $\varepsilon$ if
    \begin{enumerate}
        \item $c$ has at most one nonprivate gate;
        \item the cell of every gate of $c$ lies in $S$;
        \item every monomial of every gate expansion has at most $r$ pair effects and the
              coefficient sum of every gate expansion is at most $K$;
        \item every gate expansion has at most $K$ monomials;
        \item $\lVert\operatorname{op}(c)x-y\rVert\leq\varepsilon$.
    \end{enumerate}
    Two lists $\mathbf F$, $\mathbf F'$ of sheets with an idle system $W$ and
    $\xi\in\mathcal H_W$ have a bounded change if some original circuit from
    $\mathrm{St}(\mathbf F;W)$ to $\mathrm{St}(\mathbf F';W)$ is a bounded change from
    $\omega_{\mathbf F}(\xi)$ to $\omega_{\mathbf F'}(\xi)$, and a source-bounded change with
    constant $r_s$ if moreover that circuit satisfies the pair-source bound $r_s$
    \cite[Proposition~7.1]{OpenAI2026PolynomialPEPS}, \cite[Theorem~5.2]{OpenAI2026PolynomialPEPS}.
    A source-bounded change is a bounded change, and all three notions persist when $S$ is
    enlarged and $r$, $K$, $\varepsilon$ and $r_s$ are increased.
\end{definition}

\begin{theorem}[Bounded changes and circuits without nonprivate gates]
    \label{thm:party_bounded_change_basic}
    \lean{QuantumCircuit.Party.OriginalCircuit.norm_toEuclideanLin_op_le,
        QuantumCircuit.Party.OriginalCircuit.IsBoundedChange.comp_right,
        QuantumCircuit.Party.OriginalCircuit.IsBoundedChange.comp_left,
        QuantumCircuit.Party.OriginalCircuit.isBoundedChange_of_nonprivateCount_eq_zero,
        QuantumCircuit.Party.exists_isBoundedChange_of_eq,
        QuantumCircuit.Party.OriginalCircuit.isEmpty_nonprivateLocations_of_nonprivateCount_eq_zero,
        QuantumCircuit.Party.OriginalCircuit.forall_participants_of_nonprivateCount_eq_zero}
    \leanok
    \uses{def:party_bounded_change}
    An original circuit does not increase the norm of a vector. An original circuit without
    nonprivate gates has no gate occurrence, so every condition on the participants of its gates
    holds, and it is a bounded change from $x$ to $\operatorname{op}(c)x$ for every set $S$, all
    $r,K\in\N$ and every $\varepsilon\geq0$; in particular, for equal sheets $F=F'$ the empty
    circuit is a bounded change from $r_F$ to $r_{F'}$. If $c$ is a bounded change from $x$ to $y$
    and $e$ has no nonprivate gate, then $c$ followed by $e$ is a bounded change from $x$ to
    $\operatorname{op}(e)y$, and $e$ followed by $c$ is a bounded change from every $x'$ with
    $\operatorname{op}(e)x'=x$ to $y$, with the same $S$, $r$, $K$ and $\varepsilon$.
\end{theorem}
\begin{proof}\leanok
    \uses{thm:party_op_contraction, def:party_original_frame_sequence}
    The operator of an original circuit is a contraction by
    Theorem~\ref{thm:party_op_contraction}. A circuit without nonprivate gates has as many gate
    occurrences as nonprivate gates, none, and satisfies every bound vacuously, with error $0$.
    Under composition the numbers of nonprivate gates add and the bounds hold factor by factor;
    the errors satisfy
    $\lVert\operatorname{op}(e)(\operatorname{op}(c)x-y)\rVert\leq\lVert\operatorname{op}(c)x-y\rVert$
    and $\operatorname{op}(c)\operatorname{op}(e)x'=\operatorname{op}(c)x$.
\end{proof}

\begin{theorem}[Words and chains with one pair resource are bounded changes]
    \label{thm:party_resources_bounded}
    \lean{QuantumCircuit.Party.Word.isBoundedChange_toOriginalCircuit,
        QuantumCircuit.Party.Word.isBoundedChange_toOriginalCircuit_of_sourceCount_eq_zero,
        QuantumCircuit.Party.Monomial.isBoundedChange_toOriginalCircuit}
    \leanok
    \uses{def:party_bounded_change, def:party_word, def:party_pair_gates, def:party_allowed}
    Read an allowed word, or an allowed monomial of local contractions, normalized pair sources
    and normalized pair effects, as an original circuit whose local contractions are private and
    whose pair sources and pair effects are gates on the cells of their two parties, each with the
    single monomial ``this gate'' of coefficient one. Let $x$ and $y$ be vectors with
    $\lVert\operatorname{op}x-y\rVert\leq\varepsilon$ for the operator of the word or
    monomial.
    \begin{enumerate}
        \item An allowed word that uses only the parties of $S$ and has at most one pair source is
              a bounded change from $x$ to $y$ with participants in $S$, every $r$, every
              $K\geq1$ and error $\varepsilon$; these are the words of the births of
              Theorem~\ref{thm:party_birth}.
        \item An allowed word without pair sources is a bounded change from $x$ to $y$ for every
              $S$, $r$ and $K$, with error $\varepsilon$; these are the words of the exchanges of
              Theorem~\ref{thm:party_exchange}.
        \item An allowed monomial that uses only the parties of $S$ and has at most one gate is a
              bounded change from $x$ to $y$ with participants in $S$, every $r\geq1$, every
              $K\geq1$ and error $\varepsilon$; these are the chains of the deaths of
              Theorem~\ref{thm:party_birth}.
    \end{enumerate}
\end{theorem}
\begin{proof}\leanok
    \uses{def:party_original_circuit, def:party_word}
    The original circuit has one nonprivate gate for every pair source and pair effect, whose
    cell is the pair of its two parties, inside $S$ when the word or monomial uses only $S$. Its
    gate expansions have one monomial, of coefficient one, with at most one pair effect, so the
    expansion and monomial bounds hold for $r\geq1$ and $K\geq1$, and for all $r$ and $K$ when
    there is no gate. Its operator is that of the word or monomial.
\end{proof}

\begin{theorem}[Bounded changes of the stack of sheets]
    \label{thm:party_stack_bounded_change}
    \lean{QuantumCircuit.Party.hasSourceBoundedChange_of_perm,
        QuantumCircuit.Party.hasBoundedChange_of_perm,
        QuantumCircuit.Party.hasSourceBoundedChange_refl,
        QuantumCircuit.Party.hasBoundedChange_refl,
        QuantumCircuit.Party.HasSourceBoundedChange.perm_left,
        QuantumCircuit.Party.HasBoundedChange.perm_left,
        QuantumCircuit.Party.HasSourceBoundedChange.perm_right,
        QuantumCircuit.Party.HasBoundedChange.perm_right,
        QuantumCircuit.Party.isBoundedChange_liftOne,
        QuantumCircuit.Party.hasBoundedChange_liftOne,
        QuantumCircuit.Party.hasSourceBoundedChange_liftOne,
        QuantumCircuit.Party.isBoundedChange_liftTwo,
        QuantumCircuit.Party.hasBoundedChange_liftTwo,
        QuantumCircuit.Party.hasSourceBoundedChange_liftTwo,
        QuantumCircuit.Party.hasBoundedChange_activate,
        QuantumCircuit.Party.hasSourceBoundedChange_activate,
        QuantumCircuit.Party.hasBoundedChange_perm_liftOne,
        QuantumCircuit.Party.hasSourceBoundedChange_perm_liftOne}
    \leanok
    \uses{def:party_bounded_change, def:party_sheet_stack}
    Fix an idle system $W$ and $\xi\in\mathcal H_W$.
    \begin{enumerate}
        \item If $\mathbf F'$ is a permutation of $\mathbf F$, then $\mathbf F$, $\mathbf F'$ have
              a source-bounded change, hence a bounded change, for every $S$, $r$, $K$, $r_s$ and
              $\varepsilon\geq0$; in particular every list has the identity change. A
              (source-)bounded change from $\mathbf G$ to $\mathbf F'$ is one from every
              permutation $\mathbf F$ of $\mathbf G$, and one from $\mathbf F$ to $\mathbf G$ is one
              from $\mathbf F$ to every permutation $\mathbf F'$ of $\mathbf G$, with the same
              constants.
        \item Let $\lVert r_G\rVert\leq1$ for every sheet $G$ and $\lVert\xi\rVert\leq1$. If $c$ is
              a bounded change of $F$ into $F'$ from $r_F$ to $r_{F'}$, then $L_1(c)$ is a bounded
              change from $\omega_{(\mathbf G,F,\mathbf H)}(\xi)$ to
              $\omega_{(\mathbf G,F',\mathbf H)}(\xi)$ with the same participants, bounds and
              error, and a source-bounded one with constant $r_s$ if $c$ satisfies the pair-source
              bound $r_s$. The same holds for $L_2(c)$ when $c$ is a bounded change of $F_1,F_2$
              into $F_1',F_2'$ from $r_{F_1}\otimes r_{F_2}$ to $r_{F_1'}\otimes r_{F_2'}$.
        \item Under the norm conditions of the second item, if $\mathbf F$ is a permutation of
              $(\mathbf G,F_1,F_2,\mathbf H)$, for instance with an idle auxiliary sheet $F_2$ moved
              next to the main sheet $F_1$, a bounded change of the two adjacent sheets as in the
              second item gives a (source-)bounded change from $\mathbf F$ to
              $(\mathbf G,F_1',F_2',\mathbf H)$; likewise, if $\mathbf F$ is a permutation of
              $(\mathbf G,F,\mathbf H)$, a bounded change of $F$ gives one from $\mathbf F$ to
              $(\mathbf G,F',\mathbf H)$.
    \end{enumerate}
\end{theorem}
\begin{proof}\leanok
    \uses{thm:party_stack_perm, thm:party_bounded_change_basic, thm:party_stack_lift_error,
        def:party_sheet_stack, def:party_original_frame_sequence}
    A permutation of the sheets is the relabelling of Theorem~\ref{thm:party_stack_perm}, an
    original circuit without nonprivate gates, hence without pair sources, that maps one
    reference vector exactly to the other; Theorem~\ref{thm:party_bounded_change_basic} gives
    the first item. For the second, lifting keeps the number of nonprivate gates, the
    participants and the bounds by Definition~\ref{def:party_sheet_stack}, and
    Theorem~\ref{thm:party_stack_lift_error} bounds the error. The third item combines the second
    with the reordering of the first.
\end{proof}

\begin{theorem}[Private product initialization]
    \label{thm:party_stack_initialization}
    \lean{QuantumCircuit.Party.partyFactor, QuantumCircuit.Party.eq_of_apply_holder_eq,
        QuantumCircuit.Party.prod_partyFactor, QuantumCircuit.Party.norm_partyFactor,
        QuantumCircuit.Party.exists_productVector_eq_framesVec,
        QuantumCircuit.Party.Sheet.exists_productVector_eq_framesVec,
        QuantumCircuit.Party.exists_originalCircuit_framesVec,
        QuantumCircuit.Party.Sheet.exists_originalCircuit_framesVec}
    \leanok
    \uses{def:party_sheet_stack, def:party_sheet, def:party_original_circuit,
        def:siteconfig_product_vector}
    Let $h:\kappa\to\mathcal P$ assign a holder to every sheet, and let $\mathbf F$ be a list of
    sheets such that $X^F_p$ is a single point for every $F$ in $\mathbf F$ and every $p\neq h(F)$.
    \begin{enumerate}
        \item A configuration $z$ of $X^F$ is determined by $z_{h(F)}$. The factor of $F$ at the
              party $p$ is the vector $\pi_{F,p}$ on $X^F_p$ with
              $\pi_{F,p}(a)=\sum_{z:z_p=a}r_F(z)$ if $p=h(F)$ and $\pi_{F,p}=1$ otherwise; then
              $\prod_p\pi_{F,p}(z_p)=r_F(z)$, and $\lVert\pi_{F,p}\rVert$ is $\lVert r_F\rVert$ for
              $p=h(F)$ and $1$ otherwise.
        \item There are vectors $\varphi_p$ on $\mathrm{St}(\mathbf F;W)_p$, with $W$ the system of
              one basis vector at every party, such that
              $\lVert\varphi_p\rVert=\prod_{F\in\mathbf F,\,h(F)=p}\lVert r_F\rVert$ and the product
              vector $\bigotimes_p\varphi_p$ is $\omega_{\mathbf F}$.
        \item If moreover $\lVert r_F\rVert=1$ for every $F$ in $\mathbf F$, there is an original
              circuit without nonprivate gate from the system of one basis vector at every party
              to $\mathrm{St}(\mathbf F;W)$ whose operator maps that basis vector to
              $\omega_{\mathbf F}$.
        \item For a list of ownership sheets $S$ each owned entirely by $h(S)$, a unit vector
              $\Omega\in\C^{[q]^\Lambda}$ and $r_S=\Omega_S$, the second item holds with
              $\lVert\varphi_p\rVert=1$ for every party, and the third item holds.
    \end{enumerate}
    So the initial memory of \cite[Proposition~7.1]{OpenAI2026PolynomialPEPS}, every copy of
    $\Omega$ initialized privately at its future owner, is a pure product across the parties,
    prepared without nonprivate gates.
\end{theorem}
\begin{proof}\leanok
    \uses{def:party_sheet_stack, def:party_sheet, def:party_word,
        thm:party_resources_bounded}
    Two configurations of $X^F$ that agree at $h(F)$ agree everywhere, since the other bases are
    single points; so the sum defining $\pi_{F,h(F)}(a)$ has the single term $r_F(z)$ with
    $z_{h(F)}=a$, which gives the product formula and, the map $z\mapsto z_{h(F)}$ being a
    bijection, the norm. The second item follows by induction on the list: for $(F,\mathbf F')$
    take $\varphi_p=\pi_{F,p}\otimes\varphi'_p$, whose norm is
    $\lVert\pi_{F,p}\rVert\,\lVert\varphi'_p\rVert$ and whose product over the parties is
    $r_F\otimes\omega_{\mathbf F'}$. For the third, every $\varphi_p$ is a unit vector, and the
    word preparing $\varphi_p$ at every party by one local preparation is allowed and has no
    pair source, so as an original circuit it has no nonprivate gate, and it maps the basis
    vector to $\bigotimes_p\varphi_p$. For a sheet owned entirely by $h(S)$, every other party owns
    no site, so its basis $[q]^\emptyset$ is a single point, and
    $\lVert\Omega_S\rVert=\lVert\Omega\rVert=1$ by Definition~\ref{def:party_sheet}.
\end{proof}
